\documentclass[../../main.tex]{subfiles}
\begin{document}

\begin{center}
  {\fontsize{48pt}{58pt}\selectfont \textbf{2006}\normalsize\textbf{年度}}
\end{center}

\begin{table}[htb]
  \centering
  \caption{出題テーマと難易度}
  \begin{tabular}{|c|p{9cm}|c|}
    \hline
    \textbf{問題番号} & \textbf{テーマ} & \textbf{難易度} \\
    \hline
    第1問 & 一次式の平方和と比例関係 & \\
    \hline
    第2問 & 行列のべきと数列の収束 & \\
    \hline
    第3問 & さいころの目の和の確率 & \\
    \hline
    第4問 & 三角形の内接円と外接円 & \\
    \hline
    第5問 & 円錐状の容器の排水時間 & \\
    \hline
    第6問 & 正接の値の無理性 & \\
    \hline
  \end{tabular}
\end{table}

\begin{table}[htb]
  \centering
  \caption{解答}
  \renewcommand{\arraystretch}{1.4}
  \begin{tabular}{|c|c|p{9cm}|}
    \hline
    \textbf{問題番号} & \textbf{小問} & \textbf{答え} \\
    \hline
    第1問 & -- & 証明問題 \\
    \hline
    第2問 & -- &  \\
    \hline
    第3問 & -- & $n(n+1)(n+2)\left(\dfrac{1}{6}\right)^{n+1}$ \\
    \hline
    第4問 & -- & $r\le\dfrac{1}{2}$，等号は正三角形のとき \\
    \hline
    第5問 & -- & $\dfrac{2}{5}\pi R^2\sqrt{H}$ \\
    \hline
    第6問 & -- & $\tan 1^\circ$は無理数 \\
    \hline
  \end{tabular}
\end{table}

\paragraph{試験形態}
6問．試験時間は要確認．

\end{document}
